% !TEX program = pdflatex
\documentclass[10pt,aspectratio=169]{beamer}

\usepackage[utf8]{inputenc}
\usepackage[T1]{fontenc}
\usepackage[english]{babel}
\usepackage{lmodern}
\usepackage{microtype}
\usepackage{amsmath,amssymb,bm,mathtools}
\usepackage{booktabs,array}

\usetheme[numbering=fraction,progressbar=frametitle]{metropolis}
\setbeamertemplate{frame numbering}[fraction]
\setbeamertemplate{navigation symbols}{}

\definecolor{CourseBlue}{HTML}{1A3C68}
\definecolor{CourseGreen}{HTML}{216E39}
\definecolor{CourseGold}{HTML}{B8860B}
\definecolor{CourseRed}{HTML}{9A2C2C}
\definecolor{SoftBlue}{HTML}{F2F6FA}
\definecolor{SoftGreen}{HTML}{F1F8F3}
\definecolor{SoftRed}{HTML}{FFF3F3}

\setbeamercolor{normal text}{fg=black,bg=white}
\setbeamercolor{structure}{fg=CourseBlue}
\setbeamercolor{title}{fg=CourseBlue}
\setbeamerfont{title}{size=\Large,series=\bfseries}
\setbeamercolor{frametitle}{bg=CourseBlue,fg=white}
\setbeamerfont{frametitle}{series=\bfseries,size=\large}
\setbeamertemplate{frametitle}{
  \nointerlineskip
  \begin{beamercolorbox}[wd=\paperwidth,ht=3.2ex,dp=1.4ex,leftskip=1em,rightskip=1em]{frametitle}
    \usebeamerfont{frametitle}\insertframetitle
  \end{beamercolorbox}
}
\setbeamercolor{block title}{bg=CourseBlue,fg=white}
\setbeamercolor{block body}{bg=SoftBlue,fg=black}
\setbeamercolor{block title example}{bg=CourseGreen,fg=white}
\setbeamercolor{block body example}{bg=SoftGreen,fg=black}
\setbeamercolor{block title alerted}{bg=CourseRed,fg=white}
\setbeamercolor{block body alerted}{bg=SoftRed,fg=black}
\setbeamertemplate{blocks}[rounded][shadow=false]

\newenvironment{modelbox}{\begin{beamercolorbox}[rounded=true,sep=0.7em,wd=\textwidth]{block body}
  \textbf{\color{CourseBlue}MODEL}\quad}{\end{beamercolorbox}}
\newenvironment{marketbox}{\begin{beamercolorbox}[rounded=true,sep=0.7em,wd=\textwidth]{block body example}
  \textbf{\color{CourseGreen}NO-ARBITRAGE}\quad}{\end{beamercolorbox}}
\newenvironment{simplificationbox}{\begin{beamercolorbox}[rounded=true,sep=0.7em,wd=\textwidth]{block body alerted}
  \textbf{\color{CourseRed}WATCH OUT}\quad}{\end{beamercolorbox}}
\newenvironment{pausebox}{\begin{beamercolorbox}[rounded=true,sep=0.8em,wd=\textwidth]{block body}
  \textbf{\color{CourseGold}YOUR TURN}\quad}{\end{beamercolorbox}}

\newcommand{\E}{\mathbb E}
\newcommand{\Q}{\mathbb Q}
\newcommand{\F}{\mathcal F}
\newcommand{\dd}{\,\mathrm d}
\newcommand{\Pz}{P}

\AtBeginSection[]{
  \begin{frame}[plain]\vfill
    {\Large\bfseries\color{CourseBlue}\insertsectionhead\par}
    \medskip{\color{CourseBlue}\rule{0.92\textwidth}{0.8pt}}\vfill
  \end{frame}
}

\title{Change of Numeraire}
\subtitle{Choosing the unit that makes pricing simple}
\author{}
\date{Lecture 11}

\begin{document}

\begin{frame}[plain]
  \vfill
  {\LARGE\bfseries\color{CourseBlue}Change of Numeraire\par}
  \medskip{\large Choosing the unit that makes pricing simple\par}
  \medskip{\color{CourseGold}\rule{0.94\textwidth}{0.8pt}}
  \medskip{Lecture 11\par}\vfill
\end{frame}

\begin{frame}{One price, many useful representations}
\begin{block}{The question}
Risk-neutral pricing discounts by the money-market account. What changes if a
different positive traded asset is used as the unit of account?
\end{block}
\begin{enumerate}
  \item Define a numeraire and the measure that comes with it.
  \item Derive the pricing formula in any unit of account.
  \item See why a bond numeraire produces a forward measure.
  \item Apply the same logic to domestic and foreign currency viewpoints.
\end{enumerate}
\begin{marketbox}
The no-arbitrage price does not change. We change only its representation:
the unit, the measure, and the martingale process.
\end{marketbox}
\end{frame}

\section{Units, measures, and martingales}

\begin{frame}{Risk-neutral valuation already uses a numeraire}
With money-market account \(B_t\), the usual pricing formula is
\[
V_t=B_t\,\E^\Q\!\left[\frac{X_T}{B_T}\mid\F_t\right].
\]
Equivalently, for a constant rate \(r\),
\[
V_t=\E^\Q\!\left[e^{-r(T-t)}X_T\mid\F_t\right].
\]
\begin{block}{What is a numeraire?}
It is the strictly positive traded asset used as the unit in which prices
are quoted. Here, one unit means one unit of cash invested in \(B\).
\end{block}
\begin{marketbox}
Under the measure associated with \(B\), every traded price divided by
\(B_t\) is a martingale.
\end{marketbox}
\end{frame}

\begin{frame}{A numeraire is a change of unit, not a new economy}
Let \(N_t>0\) be a traded asset. Its denominated price for an asset \(S\) is
\[
S_t^{(N)}=\frac{S_t}{N_t}.
\]
\begin{center}
\begin{tabular}{c|c|c}
\toprule
Choice of unit & Martingale object & Associated measure\\
\midrule
Money-market account \(B_t\) & \(S_t/B_t\) & \(\Q\)\\
Zero-coupon bond \(P(t,T)\) & \(S_t/P(t,T)\) & \(\Q^T\)\\
Stock \(S_t\) & \(X_t/S_t\) & \(\Q^S\)\\
\bottomrule
\end{tabular}
\end{center}
\begin{simplificationbox}
Dividing by \(N_t\) but keeping the old measure is generally wrong. The
numeraire and probability measure must change together.
\end{simplificationbox}
\end{frame}

\begin{frame}{Ratios carry an It\^o correction}
Suppose, under one measure,
\[
\frac{\dd S_t}{S_t}=\mu_S\dd t+\sigma_S\dd W_t,\qquad
\frac{\dd N_t}{N_t}=\mu_N\dd t+\sigma_N\dd W_t.
\]
Then the ratio satisfies
\[
\frac{\dd(S_t/N_t)}{S_t/N_t}
=
\bigl(\mu_S-\mu_N+\sigma_N^2-\sigma_S\sigma_N\bigr)\dd t
+(\sigma_S-\sigma_N)\dd W_t.
\]
\begin{pausebox}
Is the drift of \(S_t/N_t\) simply \(\mu_S-\mu_N\)?
\end{pausebox}
\pause
\begin{alertblock}{No}
The quadratic-variation terms are real. The change of measure is precisely
what removes the remaining drift of the denominated price.
\end{alertblock}
\end{frame}

\begin{frame}{Every positive traded numeraire has its own martingale measure}
Start under the money-market measure \(\Q\). The measure associated with
a positive traded numeraire \(N\) is defined by
\[
\boxed{
\left.\frac{\dd\Q^N}{\dd\Q}\right|_{\F_t}
=
\frac{N_t/B_t}{N_0/B_0}.}
\]
\begin{block}{Why this is a valid density}
\(N_t/B_t\) is a positive \(\Q\)-martingale because \(N\) is traded.
After normalizing at time \(0\), its expectation is one.
\end{block}
\begin{marketbox}
Under \(\Q^N\), every traded asset price divided by \(N_t\) is a martingale.
\end{marketbox}
\end{frame}

\begin{frame}{Changing measure reweights states; it does not create states}
For two positive numeraires \(N\) and \(M\),
\[
\boxed{
\left.\frac{\dd\Q^M}{\dd\Q^N}\right|_{\F_t}
=
\frac{M_t/N_t}{M_0/N_0}.}
\]
\begin{columns}[T,onlytextwidth]
\column{0.48\textwidth}
\begin{block}{Economic reading}
States in which \(M\) performs well relative to \(N\) receive more weight
under \(\Q^M\) than under \(\Q^N\).
\end{block}
\column{0.48\textwidth}
\begin{exampleblock}{What stays fixed}
The measures are equivalent: they assign zero probability to exactly the
same events.
\end{exampleblock}
\end{columns}
\begin{pausebox}
Does a change of numeraire alter which market scenarios are possible?
\end{pausebox}
\pause
\begin{exampleblock}{Answer}
No. It changes probability weights, not the scenarios themselves.
\end{exampleblock}
\end{frame}

\begin{frame}{The general pricing formula has the same three steps}
For a claim paying \(X_T\) at time \(T\),
\[
\boxed{
V_t=N_t\,\E^{\Q^N}\!\left[\frac{X_T}{N_T}\mid\F_t\right].}
\]
\begin{enumerate}
  \item Divide the payoff by the numeraire at the payment date.
  \item Take conditional expectation under the associated measure.
  \item Multiply by the current numeraire value.
\end{enumerate}
\begin{marketbox}
A good numeraire makes the normalized payoff \(X_T/N_T\) simple and makes
the relevant normalized price a martingale.
\end{marketbox}
\end{frame}

\begin{frame}{The money-market formula is one special case}
Set \(N_t=B_t\). Then \(\Q^N=\Q\) and
\[
V_t=B_t\,\E^\Q\!\left[\frac{X_T}{B_T}\mid\F_t\right].
\]
For deterministic \(r\),
\[
\frac{B_t}{B_T}=e^{-r(T-t)},
\]
which gives the familiar discounted expectation.
\begin{pausebox}
What is the density process when the new numeraire is again \(B_t\)?
\end{pausebox}
\pause
\begin{exampleblock}{Answer}
It is identically one. Choosing the money-market account does not change
the usual risk-neutral measure.
\end{exampleblock}
\end{frame}

\section{Forward measures simplify payment-date pricing}

\begin{frame}{A zero-coupon bond absorbs discounting}
Fix payment date \(T\) and choose the bond \(P(t,T)\) as numeraire. Since
\(P(T,T)=1\),
\[
\boxed{
V_t=P(t,T)\,\E^{\Q^T}[X_T\mid\F_t].}
\]
\begin{columns}[T,onlytextwidth]
\column{0.48\textwidth}
\begin{block}{Under \(\Q\)}
Discount factor and payoff may both be random:
\[
V_t=\E^\Q\!\left[e^{-\int_t^T r_s\dd s}X_T\mid\F_t\right].
\]
\end{block}
\column{0.48\textwidth}
\begin{exampleblock}{Under \(\Q^T\)}
The discounting is embedded in \(P(t,T)\); only the payoff remains in the
expectation.
\end{exampleblock}
\end{columns}
\end{frame}

\begin{frame}{Forward prices are martingales under their payment-date measure}
For an asset with no intermediate cash flows, define
\[
F(t,T)=\frac{S_t}{P(t,T)}.
\]
Under the \(T\)-forward measure,
\[
\boxed{F(t,T)=\E^{\Q^T}[S_T\mid\F_t].}
\]
\begin{block}{What this buys us}
The forward price has zero drift under the measure naturally attached to the
delivery date. This is often much cleaner than tracking a stochastic short
rate under the money-market measure.
\end{block}
\end{frame}

\begin{frame}{A forward contract is priced by one martingale condition}
A forward on a non-dividend-paying asset delivers one share at \(T\) for
delivery price \(K\). Its value is
\[
V_t=S_t-KP(t,T)=P(t,T)\bigl(F(t,T)-K\bigr).
\]
\begin{pausebox}
If \(S_0=100\) and \(P(0,T)=0.95\), what delivery price makes the forward
worth zero at inception?
\end{pausebox}
\pause
\begin{exampleblock}{Answer}
Set \(K=F(0,T)=S_0/P(0,T)=105.26\). The forward measure did not change the
economic price; it exposed the natural martingale behind it.
\end{exampleblock}
\end{frame}

\begin{frame}{Choose the numeraire that cancels a payoff factor}
Suppose a claim has payoff \(X_T=S_TY_T\). Under the stock numeraire,
\[
\boxed{V_t=S_t\,\E^{\Q^S}[Y_T\mid\F_t].}
\]
The factor \(S_T\) has disappeared from the expectation.
\begin{columns}[T,onlytextwidth]
\column{0.48\textwidth}
\begin{block}{Good use case}
Payoffs naturally proportional to the terminal stock, such as exchange
options or share-settled claims.
\end{block}
\column{0.48\textwidth}
\begin{alertblock}{Condition}
The candidate must be a strictly positive traded asset. A payoff that can
hit zero is not automatically a valid numeraire.
\end{alertblock}
\end{columns}
\end{frame}

\begin{frame}{A practical numeraire-selection checklist}
\begin{enumerate}
  \item What is the payoff date and currency?
  \item Does the payoff contain a positive traded factor?
  \item Does a bond maturing on the payment date remove discounting?
  \item Under which measure is the relevant ratio a martingale?
  \item Can the final answer be checked against money-market pricing?
\end{enumerate}
\begin{pausebox}
A payoff is paid at \(T\) and depends on future bond prices. Which first
candidate for the numeraire should you test?
\end{pausebox}
\pause
\begin{exampleblock}{Answer}
The zero-coupon bond \(P(t,T)\). It matches the payment date and generates
the \(T\)-forward measure.
\end{exampleblock}
\end{frame}

\section{FX makes the two viewpoints concrete}

\begin{frame}{An exchange rate links two money-market accounts}
Let \(X_t\) be domestic currency per unit of foreign currency. With constant
rates,
\[
B_d(t)=e^{r_dt},\qquad B_f(t)=e^{r_ft}.
\]
\begin{center}
\begin{tabular}{c|c|c}
\toprule
Object & Units & Domestic value\\
\midrule
One foreign currency unit & foreign & \(X_t\) domestic\\
Foreign bank account & foreign & \(\boxed{X_tB_f(t)}\) domestic\\
Domestic bank account & domestic & \(B_d(t)\) domestic\\
\bottomrule
\end{tabular}
\end{center}
\begin{marketbox}
\(X_tB_f(t)\) is a positive traded asset in domestic currency. It is the
natural bridge between domestic and foreign pricing viewpoints.
\end{marketbox}
\end{frame}

\begin{frame}{Domestic no-arbitrage fixes the FX drift}
Under the domestic money-market numeraire, the discounted domestic value of
foreign cash must be a martingale:
\[
\frac{X_tB_f(t)}{B_d(t)}
\quad\text{is a }\Q^d\text{-martingale}.
\]
Therefore
\[
\boxed{
\frac{\dd X_t}{X_t}=(r_d-r_f)\dd t+\sigma\dd W_t^{\Q^d}.}
\]
\begin{block}{Interpretation}
The foreign rate acts exactly like a continuous dividend yield. Holding
foreign currency gives access to foreign interest, just as holding a stock
can generate dividends.
\end{block}
\end{frame}

\begin{frame}{The FX forward is the first diagnostic}
For delivery of one unit of foreign currency at \(T\),
\[
\boxed{F_{0,T}=X_0e^{(r_d-r_f)T}.}
\]
\begin{pausebox}
Take \(X_0=1.10\) USD/EUR, \(r_d=5\%\), \(r_f=3\%\), and \(T=1\).
What is the no-arbitrage one-year forward rate in USD/EUR?
\end{pausebox}
\pause
\begin{exampleblock}{Answer}
\[
F_{0,1}=1.10e^{0.02}=1.1222\ \text{USD/EUR}.
\]
Domestic interest pushes the forward up; foreign interest pushes it down.
\end{exampleblock}
\end{frame}

\begin{frame}{An FX call is Black--Scholes with a foreign yield}
A domestic-currency call on one unit of foreign currency pays
\((X_T-K)^+\). Its domestic value is
\[
\boxed{
C_t=X_te^{-r_f\tau}\Phi(d_1)-Ke^{-r_d\tau}\Phi(d_2),}
\]
where \(\tau=T-t\) and
\[
d_1=\frac{\ln(X_t/K)+(r_d-r_f+\tfrac12\sigma^2)\tau}
{\sigma\sqrt{\tau}},\qquad d_2=d_1-\sigma\sqrt{\tau}.
\]
\begin{marketbox}
Read this as the usual dividend-paying-stock formula: domestic rate discounts
the strike, foreign rate discounts the spot.
\end{marketbox}
\end{frame}

\begin{frame}{Worked FX option: identify the two discount factors}
Let \(X_0=1.10\), \(K=1.12\), \(r_d=5\%\), \(r_f=3\%\),
\(\sigma=12\%\), and \(T=1\). Then
\[
d_1=0.077,\qquad d_2=-0.043,
\]
and
\[
C_0=1.10e^{-0.03}\Phi(0.077)
-1.12e^{-0.05}\Phi(-0.043)\approx0.052.
\]
\begin{block}{Reading the answer}
The option costs approximately \(0.052\) USD per EUR of notional. The two
rates appear in different places because the contract exchanges two currencies.
\end{block}
\end{frame}

\begin{frame}{FX put--call parity reveals the same structure}
For domestic prices \(C_t\) and \(P_t\) with the same \(K,T\),
\[
\boxed{
C_t-P_t=X_te^{-r_f(T-t)}-Ke^{-r_d(T-t)}.}
\]
\begin{pausebox}
Which term is the present value of one unit of foreign currency received at
\(T\): \(X_te^{-r_f(T-t)}\) or \(X_te^{-r_d(T-t)}\)?
\end{pausebox}
\pause
\begin{exampleblock}{Answer}
\(X_te^{-r_f(T-t)}\). Invest one foreign unit at the foreign rate, then
convert its future value into domestic currency.
\end{exampleblock}
\end{frame}

\begin{frame}{Domestic and foreign measures are linked by the foreign bank account}
In domestic currency, choose
\[
N_t^f=X_tB_f(t)
\]
as the foreign-cash numeraire. The density from the domestic measure is
\[
\boxed{
\left.\frac{\dd\Q^f}{\dd\Q^d}\right|_{\F_t}
=
\frac{X_tB_f(t)/B_d(t)}{X_0B_f(0)/B_d(0)}.}
\]
\begin{block}{The two viewpoints}
Domestic pricing measures values in domestic cash; foreign pricing measures
values relative to foreign cash. They give the same no-arbitrage price after
the appropriate currency conversion.
\end{block}
\end{frame}

\begin{frame}{A change of numeraire is useful only with the matching measure}
\begin{center}
\begin{tabular}{p{0.30\textwidth}|p{0.30\textwidth}|p{0.25\textwidth}}
\toprule
Pricing situation & Helpful numeraire & Martingale quantity\\
\midrule
Ordinary equity claim & Money-market account & \(S_t/B_t\)\\
Payoff at one fixed date & Bond \(P(t,T)\) & \(S_t/P(t,T)\)\\
Payoff proportional to stock & Stock \(S_t\) & \(X_t/S_t\)\\
FX in domestic currency & Domestic cash \(B_d\) & \(X_tB_f(t)/B_d(t)\)\\
\bottomrule
\end{tabular}
\end{center}
\begin{simplificationbox}
Never mix the numerator from one convention with the measure from another.
That is the most common source of sign and drift mistakes.
\end{simplificationbox}
\end{frame}

\section{Takeaways}

\begin{frame}{The lecture reduces to one reusable formula}
\[
\boxed{
V_t=N_t\,\E^{\Q^N}\!\left[\frac{X_T}{N_T}\mid\F_t\right].}
\]
\begin{enumerate}
  \item A numeraire is a positive traded unit of account.
  \item It determines an equivalent measure under which denominated prices
  are martingales.
  \item The price is invariant; the representation can become much simpler.
  \item Bond numeraires create forward measures; FX supplies a natural
  two-currency application.
\end{enumerate}
\end{frame}

\begin{frame}{Final check: choose the unit before doing the algebra}
\begin{pausebox}
A claim pays at date \(T\) in domestic currency and its payoff is a function
of a \(T\)-maturity forward rate. What numeraire is the natural first choice,
and what is the resulting pricing measure?
\end{pausebox}
\pause
\begin{exampleblock}{Answer}
Start with the zero-coupon bond \(P(t,T)\). It leads to the \(T\)-forward
measure and removes discounting from the expectation at the payment date.
\end{exampleblock}
\vfill
\begin{center}
{\large\color{CourseBlue}\bfseries The numeraire changes the view, not the value.}
\end{center}
\end{frame}

\end{document}
